\section{The objective and its derivatives}
\label{sec:lvcalib}

\subsection{The objective}

The pieces now assemble into a bound-constrained least-squares problem
over the nodal variances:
\begin{equation}
\min_{\theta\in[v_{\rm lo},\,v_{\rm hi}]^{n_\theta}}
\;\tfrac12\Big(
\sum_{q}r_q(\theta)^{2}
\;+\;\lambda\,\big\|\Gamma(\theta-\theta_{\rm ref})\big\|^{2}
\;+\;\lambda_0\,\big\|\theta_{1,\cdot}-\theta_{2,\cdot}\big\|^{2}
\Big).
\label{eq:lvobjective}
\end{equation}
Its three blocks, in order: a data term, a roughness penalty, and a
tie between the first two vertex rows --- $\theta_{1,\cdot}$ is the
$\tau=0$ row, $\theta_{2,\cdot}$ the row at the first quoted expiry.
Each block was argued for before it was written down; we
take them in turn.

The \emph{data block} compares the marched call $c_\theta$ at quote
$q$'s coordinates with the quoted one, in volatility units:
$r_q=\big(c_\theta(\tau_q,y_q)-c_q\big)/\eta_q$, where $\eta_q$ is one
volatility percent of the quote's vega, floored.  Recall what vega is:
the derivative of the Black price with respect to implied volatility
--- the exchange rate between price units and volatility units.
Dividing each price error by it makes every residual read as a
volatility error, so a deep wing's tiny price errors and the money's
large ones compete on the same scale; the floor (the same device as
Chapter~2's vega floor $\eta$; this chapter's constant is in
\cref{app:lvnumerics}) keeps the scale from vanishing where vega does.
This is the residual design of \cref{sec:calibration} --- solve in
price space, measure in volatility space --- carried over unchanged,
and when band edges rather than mids are trusted, the band objective
of that section --- residuals measured to the nearest bid--ask band
edge, zero inside the band --- applies verbatim.

The \emph{roughness block} $\lambda\|\Gamma(\theta-\theta_{\rm ref})\|^{2}$
is Tikhonov regularization --- the standard name for adding a penalty
so that, among sheets fitting the data equally well, the optimizer
returns the smoothest.  The operator $\Gamma$ takes second divided
differences along strike within each maturity row and along time
within each strike column --- the discrete curvature of the sheet ---
normalized by the local vertex spacing, so that a tightly spaced money
column is not penalized more than a wide wing one (on a uniform grid
it reduces to the $(1,-2,1)$ stencil).  The penalty pulls toward a
reference sheet $\theta_{\rm ref}$ --- flat, at the median ATM
variance, in this chapter's protocol --- and it is the declared answer
to \cref{sec:lvinfo}: of the equivalence class the data allows, return
the smoothest member near the reference.

The \emph{front tie}, at weight $\lambda_0$, pins the $\tau=0$ row to
the first quoted row, closing the invisible sub-front subspace of
\cref{sec:lvinfo}.

Two further row types plug into the same stack when their instruments
are present, and are only named here.  A variance-swap quote adds one
scalar residual per quoted expiry, constraining the deep-put tail
integral that vanillas cannot see --- Chapter~5 develops the integral,
and it is exactly the instrument that justifies releasing the wing
slope $a$ of \cref{rem:lvhull}.  And prior structures from a previous
fit enter as reference terms in the same
$\Gamma$-and-$\theta_{\rm ref}$ position --- Chapter~10 develops when
and how much to trust them.

\subsection{Every derivative from the same factorization}

A least-squares solver runs on the Jacobian: the matrix of derivatives
of every residual with respect to every parameter.  For a model priced
by a PDE march, the naive route is finite differences --- perturb
$\theta_\ell$, re-march the surface, subtract --- which costs
$n_\theta$ extra marches per iteration, hundreds here.  The sheet's
structure gives something better: \emph{exact} derivatives, for the
price of extra right-hand sides through a factorization already
computed.  The key is linearity.  By \eqref{eq:lvp1} the field depends
linearly on $\theta$, so the discrete generator does too, and its
derivative in one parameter is itself a stencil:
\[
\frac{\partial\mathcal L}{\partial\theta_\ell}
=\mathcal L[\mathcal H_\ell],
\]
the stencil \eqref{eq:lvstencil} with the field $v$ replaced by the
single hat function $\mathcal H_\ell$ --- zero outside vertex $\ell$'s
star of triangles.

\begin{proposition}[Tangent: same factor, extra right-hand sides]
\label{prop:lvtangent}
The sensitivities $s^{n}_\ell=\partial U^{n}/\partial\theta_\ell$ obey
\begin{equation}
\mathcal M^{n+1}s^{n+1}_\ell
=s^{n}_\ell+\Delta\tau\,\mathcal L[\mathcal H_\ell]\,U^{n+1},
\qquad s^{0}_\ell=0,
\label{eq:lvtangent}
\end{equation}
where $\mathcal L[\mathcal H_\ell]\,U^{n+1}$ applies the boundary-inclusive
stencil to the solution extended by its fixed boundary values.  This is the
discretization of the continuous sensitivity equation
$\partial_\tau s_\ell=\tfrac12 v y^{2}\partial_{yy}s_\ell
+\tfrac12\mathcal H_\ell\,y^{2}\partial_{yy}c$.
\end{proposition}

\begin{proof}
Write one step over the interior unknowns as
$\mathcal M^{n+1}U^{n+1}=U^{n}+b^{n+1}$, with $b^{n+1}$ the boundary
vector displayed after \eqref{eq:lvstep} --- it carries the
Dirichlet values through the stencil's boundary columns, and
therefore depends on $\theta$ through the field at boundary-adjacent
nodes.  Differentiate the whole equation with respect to
$\theta_\ell$, using the product rule on the left:
\[
\Big(\partial_{\theta_\ell}\mathcal M^{n+1}\Big)U^{n+1}
+\mathcal M^{n+1}s^{n+1}_\ell
=s^{n}_\ell+\partial_{\theta_\ell}b^{n+1}.
\]
By linearity,
$\partial_{\theta_\ell}\mathcal M^{n+1}
=-\Delta\tau\,\mathcal L[\mathcal H_\ell]$ exactly --- no
approximation is involved.  Move it to the right-hand side:
\[
\mathcal M^{n+1}s^{n+1}_\ell
=s^{n}_\ell
+\Delta\tau\,\mathcal L[\mathcal H_\ell]\,U^{n+1}
+\partial_{\theta_\ell}b^{n+1},
\]
and the last two terms combine into the boundary-inclusive product of
the statement, because the boundary \emph{values} themselves ($1$ and
$0$) do not move with $\theta$ --- only the stencil weights through
which they enter do.  The start is $s^{0}_\ell=0$: the payoff does not
depend on the field.
\end{proof}

Read \eqref{eq:lvtangent} twice.  As a computation: the sensitivity
system has the \emph{same} matrix $\mathcal M^{n+1}$ as the value
system --- only the right-hand side differs --- so every sensitivity
marches behind the same tridiagonal factorization as the value solve.
The entire $n_\theta$-column Jacobian costs one factorization plus
$n_\theta$ extra back-substitutions per step; the columns form a block
of right-hand sides that a solver processes together.  As probability:
the source term says that vertex $\ell$ influences a price exactly
through the diffusion's visits to $\ell$'s star of triangles, weighted
by the convexity $\partial_{yy}c$ it finds there.

\begin{figure}[htbp]
\centering
\paperfig{\textwidth}{fig_lv_influence.pdf}
\caption{The tangent system, seen and audited.  (a)~Sensitivity of the
priced call curve to the single vertex at $(\tau,y)=(0.25,1)$, across
the four quoted expiries.  (b)~Every analytic column of the Jacobian
against a central finite difference of the full march, over all
\MacLvInflVtx\ vertices: worst relative disagreement
\MacLvInflMaxRel.}
\label{fig:lvinfluence}
\end{figure}

\Cref{fig:lvinfluence} shows both readings on the synthetic market.
In panel~(a), the sensitivity of the priced call curve to the single
vertex at $(\tau,y)=(0.25,1)$, plotted across the four quoted
expiries.  At $\tau=0.10$ the curve is identically zero: an earlier
expiry cannot feel a later vertex, because the march has not reached
it yet.  At $\tau=0.25$, the vertex's own maturity, the sensitivity is
a sharp cone centered on its strike.  At the two later expiries the
cone spreads and flattens: the diffusion gradually forgets where it
collected its variance.  In panel~(b), the audit: every analytic
Jacobian column is compared against a central finite difference of the
full march, over all \MacLvInflVtx\ vertices, and the worst relative
disagreement is \MacLvInflMaxRel.  The derivative is exact, not
approximate.

The optimization itself is then standard machinery used plainly: a
trust-region reflective least-squares iteration on
\eqref{eq:lvobjective} --- a method of the Gauss--Newton family
(linearize the residuals at the current iterate, solve the resulting
linear least squares, repeat) that handles
box bounds natively, shrinking its steps rather than leaving the box
\citep{BranchColemanLi1999} --- with the analytic Jacobian of
\cref{prop:lvtangent} supplied at every iterate and a flat seed at the
reference sheet.  The tangent route's Jacobian cost grows linearly
with $n_\theta$; the classical \emph{adjoint} alternative --- run the
march backward once through its transposed steps, collecting the whole
gradient in a single sweep, at a cost independent of the
parameter count --- is derived in \cref{app:lvnumerics} for the reader
who will need it at larger grids.
